\section{Emergence and Bottom-up Control}
\label{sec:emergence}

Section~\ref{sec:modularity} closed with a question: given that body and configuration absorb part of the work a controller would otherwise need to do, what should the controller itself do? This section examines one influential answer. Complex global behavior can arise from many simple local interactions, without an explicit top-down specification of the global outcome. The approach is known as \textit{emergence}, and it has roots both in theoretical computer science and in biology. Modularity is its natural site in robotics, because a modular body is already structured as many local units whose collective behavior is what locomotion ultimately is.

\subsection{Concept}
\label{subsec:mrr-Concept}
 
Emergence, in the precise sense used here, refers to a relationship between the local rules followed by individual elements and the global behavior of their collective. The global behavior is predictable from but not reducible to those rules; it cannot be obtained by simply summing the agents' individual states. Two of the most influential demonstrations of this idea are computational. Cellular automata and flocking models show that minimal local update rules per agent, three simple steering vectors in flocking and bit-level transitions in cellular automata, produce richly structured global patterns that no agent has been individually told to perform \cite{StephenWolframNew}. Self-organization in biological systems, from Turing's reaction-diffusion morphogenesis to the foraging trails of slime mold, points in the same direction: organized structure can arise without a designer specifying it.

In robotics, the key implication of this principle is that it relaxes the controller's task. Rather than computing a global gait and decomposing it into module-level commands, the designer specifies only a small set of rules per module and lets the global behavior arise from how those rules interact through the modules' physical connections. Cademartiri and Bishop's review of programmable self-assembly captures the design implication directly: by encoding information into the building blocks themselves, one can program the assembly process locally and obtain the desired global structure as a consequence \cite{cademartiriProgrammableSelfassembly2015}. The same logic applies to a modular robot in motion. Simple repeated unit behaviors (sliding, inflation, contraction, vibration) combined through specific connection schemes and constraints (inflated-truss topology, flexible wire links, magnetic connectors, triangulated arrangements) are orchestrated into whole-body deformation and locomotion. Several systems already encountered in this chapter make this local-to-global relationship explicit, and the next subsection reads them through the lens of emergent control.

\subsection{Related Works}
\label{subsec:emergence-related}

A first thread demonstrates emergent locomotion at its most minimal. Li et al.'s particle robotics \cite{liParticleRoboticsBased2019}, encountered in Section~\ref{sec:modularity}, lies at this end of the spectrum: each particle does only volumetric oscillation phase-modulated by a global signal, and the aggregate of many such particles exhibits robust locomotion, object transport, and phototaxis (\textit{\cref{fig:particle_robotics}}). The high-level behavior is genuinely emergent, in the sense that no particle has been individually instructed to move toward the light. Mousa et al.'s fluidic units \cite{mousaMultifunctionalFluidicUnits}, encountered in Section~\ref{sec:soft-robotics}, illustrate the same principle in a different substrate: when multiple oscillating units are mechanically coupled through a shared elastic body, they self-synchronize into a Kuramoto-style phase pattern, with no controller specifying the synchronization. Chin et al.'s AuxBot \cite{chinFlipperStyleLocomotionStrong2023} shows the principle at the level of explicit module coordination, where each module is mechanically simple and capable only of expansion or contraction, but coordinated phase-shifting between modules produces flipper-style locomotion of the whole assembly. Across these three, a common pattern recurs: a global motion that the controller never specifies arises from how local module behaviors interact through physical or topological coupling.

A second thread treats the local-rule design as a learning problem. Yu, Matthews, and colleagues' autonomous modular legs \cite{AgileLeggedLocomotion} use reinforcement learning to discover per-module controllers that, combined across many modules, yield agile locomotion across diverse body configurations (\textit{\cref{fig:agile_legged}}); the body design space and the controller space are co-optimized, so the global behavior emerges from learned local policies rather than from a hand-designed coordinator. Bio-inspired neural approaches push in a similar direction. Zhang et al.'s distributed neural locomotion controller \cite{zhangBioInspiredDistributedNeural} encodes a network of small neural circuits per joint or limb, producing robust locomotion and emergent gait adaptations from local sensorimotor feedback rather than from a centralized planner. Together, these works show that emergence is a viable control paradigm at multiple levels of sophistication, from constant-amplitude oscillation to learned per-module policies.

\subsection{Open Questions}
\label{subsec:emergence-gap}

Where the previous three sections identified specific gaps in established literatures, emergence as a control paradigm is itself well-established. What is less developed is its application to the particular research site this thesis works in: \textit{soft, modular, pneumatic} robots whose modules are physically coupled through their inflated surfaces. Most existing emergent-control demonstrations operate on rigid agents whose coupling is informational, through radio, magnets, or wires; the few that exploit material coupling, such as Mousa's fluidic units, do so in two-dimensional fluidic networks rather than three-dimensional polyhedral configurations. Whether and how local rules over inflation states can produce coordinated locomotion when the modules are passively coupled through their shared geometry is, to the author's knowledge, an open question.

Two further dimensions remain unexplored even in the broader emergence literature. The first is configuration-dependence: when the same local rules are implemented across modules arranged in different polyhedral topologies, do qualitatively different emergent gaits arise, or do similar topologies produce similar gaits? The second is the division of labor between body and controller: how much of the gait can be left to the body and material to produce, and how much must still be specified by the local rule itself? These questions sit at the boundary between the configuration-comparison gap of Section~\ref{sec:shape-changing} and the soft-modular-pneumatic gap of Section~\ref{sec:modularity}. The thesis takes them up directly through experiment: by implementing simple cellular-automaton-style local rules on inflatable modules placed at the vertices of different cospherical polyhedra, and observing what global motions, if any, arise.
